% =============================================================================
% 041 / 068
% Status: raw
% =============================================================================
% Notes:
%
% Source question:
% 你們編撰的吳語辭典，只需要先蒐集1000-3000生活常用字，就足夠使用用了吧？但編著辭典困難的點不在於詞彙數量，而是為了尊重吳語底下各方言，避免全吳語區只接受某人群、某地域的用法，因此要儘可能蒐集、整理出“數個吳語底下方言”的詞彙、發音和用法。為嚴謹起見，還要考據詞彙的歷史變化，沒辦法只靠單純蒐集現在的詞彙、發音，就編成一部辭典。這意味著，乍看之下編撰常用3000個詞彙，實質上要理清30000個沒有書面紀錄、散落在田野的詞彙？
% =============================================================================

\chapter[你們編撰的吳語辭典，只需要先蒐集1000-3000生活常用字，就足夠使用用了吧？但編著辭典困難的點不在於詞彙數量，而是為了尊重吳語底下各方言，避免全吳語區只接受某人群、某地域的用法，因此要儘可能蒐集、整理出“數個吳語底下方言”的詞彙、發音和用法。為嚴謹起見，還要考據詞彙的歷史變化，沒辦法只靠單純蒐集現在的詞彙、發音，就編成一部辭典。這意味著，乍看之下編撰常用3000個詞彙，實質上要理清30000個沒有書面紀錄、散落在田野的詞彙？]{你們編撰的吳語辭典，只需要先蒐集1000-3000生活常用字，就足夠使用用了吧？但編著辭典困難的點不在於詞彙數量，而是為了尊重吳語底下各方言，避免全吳語區只接受某人群、某地域的用法，因此要儘可能蒐集、整理出“數個吳語底下方言”的詞彙、發音和用法。為嚴謹起見，還要考據詞彙的歷史變化，沒辦法只靠單純蒐集現在的詞彙、發音，就編成一部辭典。這意味著，乍看之下編撰常用3000個詞彙，實質上要理清30000個沒有書面紀錄、散落在田野的詞彙？}
\qaanswer{
需要一再反復強調：吳語不是字中心思維，是詞中心思維。所以並不是著出常用吳語漢字這麼簡單的事。編輯最基礎的漢吳辭典是為了能讓吳語區的人在學習吳語通用語(吳語播音體語法)語法後能夠用吳語書寫文章。所以常用詞彙肯定不止3000個。從語法和諧的角度講，有些吳語原本沒有的現代概念的詞彙不能直接使用漢語普通話詞彙，因此還需要按吳語語法構詞或半構詞。既有詞彙主要是收集和甄選全吳越地區的包括徽語區詞彙。目前雖然有不少所謂的吳語辭典，但都是吳漢辭典，實際上對於使用者而言是無法當作工具書查詢吳語詞彙並用來書寫吳語文章的。所以編輯吳漢辭典的意義很大。
}

